%!TEX program = xelatex

% \documentclass{article}
\documentclass[a4paper]{article}
\usepackage[margin=25mm]{geometry}
\usepackage[utf8]{inputenc}
\usepackage[T1]{fontenc}    % use 8-bit T1 fonts

\usepackage[UTF8]{ctex}
\usepackage{indentfirst}
\usepackage{lmodern}% http://ctan.org/pkg/lm
\usepackage{setspace}
\usepackage{verbatim}
\usepackage{amsmath}
\usepackage{amsthm}
\usepackage{graphicx}
\usepackage{subfigure}
\usepackage{booktabs}
\usepackage{tabularx}
\usepackage{multirow}
\usepackage{multicol}
\usepackage{listings}
\usepackage{relsize}
\usepackage[usenames]{color}   
\usepackage[table]{xcolor}
\usepackage{array}
\usepackage[rightcaption]{sidecap}
\usepackage{hyperref}
\usepackage{float}
\usepackage{url}
\usepackage{graphicx}

% 自己画图
\usepackage{tikz}
\usepackage{mathpazo}
\usepackage{pgfplots}

% 插入动画
\usepackage{animate}
\usepackage{caption}
\usepackage{subfigure}


% 解决目录红框问题
% \hypersetup{
% colorlinks=true,
% linkcolor=black
% }
\hypersetup{
colorlinks=true,       % false: boxed links; true: colored links
linkcolor=red,          % color of internal links
citecolor=blue,        % color of links to bibliography
filecolor=magenta,      % color of file links
urlcolor=cyan   
}

% 代码格式设置
\lstset{
         language = java, numbers=left, 
         numberstyle=\tiny,keywordstyle=\color{blue!70},
         commentstyle=\color{red!50!green!50!blue!50},frame=shadowbox,
         rulesepcolor=\color{red!20!green!20!blue!20},basicstyle=\ttfamily,
         breaklines=true,               % 自动换行
         % breakatwhitespace=true,        % 只在空格处断行（更美观）
         showstringspaces=false,
         tabsize=4,
         frame=single                   % 加边框
}




\newcommand{\myd}{\;\mathrm{d}} %自定义积分变量命令
\renewcommand{\listtablename}{Tables} %改变listoftable的名字
\renewcommand{\arraystretch}{1.8} %行间距相对原始间距的倍数
% \setmonofont{Consolas}
\renewcommand\refname{参考网站}
% \setlength{\tabcolsep}{18pt} % 文本到单元格左右边框之间的距离
% \setlength{\arrayrulewidth}{0.4mm} % 设置表格线宽度


\setlength{\parindent}{2em} %控制首行缩进
\addtolength{\parskip}{3pt} %\parskip 控制段落(paragraph)距离
%\singlespacing % 设置单倍行距
%\doublespacing % 设置双倍行距
%\setstretch{1.25} % 设置行距为1.25
\onehalfspacing % 1.5倍行距
\graphicspath{{../pictures/}} % 指定图片所在文件夹
%\newtheorem{definition}[定义]{section}
%定制环境: {环境名}[编号延续]{显示名}[编号层次]

%\begin{doublespacing} 特定行距环境 \end{doublespacing}

% 让编号和小标题目录深度支持到第四级 (paragraph)
\setcounter{secnumdepth}{5}
\setcounter{tocdepth}{5}

\begin{document}
\begin{titlepage}
% titlepage格式可以自己调
        % \vspace*{-3cm}
	
	\begin{figure}[h]
		\centering
		
\includegraphics[width=0.5\linewidth]{zjdx}
	\end{figure}

	\vspace*{0.5cm}
	\begin{figure}[h]
		\centering
		
\includegraphics[width=0.3\linewidth]{QSY}
	\end{figure}
	\vspace{-0.5cm}
	\begin{center}
		\Huge{\textbf{C++项目管理和工程实践}}\\
		
		\Huge{\textbf{验收总结报告}}
	\end{center}
	
	\vspace*{0.5cm}


	\vspace*{1cm}
    \begin{center}
            \Large 小组名称\ \ \underline{\makebox[220pt]{双截龙组}} \\ 
            \vspace{0.3cm}
            \quad \Large 成员 \ \quad \underline{\makebox[220pt]{苏易文，邱安琪 3240105004，陈棋隆}} \\ 
            \vspace{0.3cm}
            \Large 实验日期\ \ \underline{\makebox[220pt]{2026年 07 月 17 日}}\\
            \vspace{0.3cm}
            \Large 指导老师\ \ \underline{\makebox[220pt]{袁昕}}\\
            

             
    \end{center}
        
    
\end{titlepage}

\newpage
% \setcounter{tocdepth}{5}

\clearpage \thispagestyle{empty} %清楚图片格式（取消掉编号）

% 不想要可以注释掉
\tableofcontents % 生成目录
\thispagestyle{empty} 

\newpage
\pagenumbering{arabic}%重新开始标号,阿拉伯数字形式

\section{成果展示}

\subsection{项目简介}
本实验仓库实现的是一款横版街机动作游戏，主题为“双截龙”。项目使用 C++17、Qt6、CMake 与 vcpkg 组织工程，当前代码已经形成了较清晰的 MVVM 分层：`common` 层提供跨层公共状态和通知工具，`viewmodel` 层负责游戏规则和状态推进，`view` 层负责 Qt 窗口、输入和绘制，`app` 层作为组合根管理生命周期以及属性、命令、通知三类绑定。

从玩法目标看，玩家扮演一名格斗角色，在带有纵深的城市街道上从左向右推进。当前实现支持标题画面、移动、跳跃、轻攻击、重攻击、血量、精力、疲劳提示、普通敌人追击、Boss 触发、暂停、失败、胜利和重开等核心流程。普通敌人在初始阶段生成并主动靠近玩家；当玩家推进到关卡后段时，系统生成 Boss，击败 Boss 后进入胜利结算。
\subsubsection{组内分工}

本组以层次边界作为分工单位，每位成员既负责对应代码，也负责对应分报告：

- 苏易文：负责 App/Common 层，主要内容包括生命周期、组合根、公共状态、通知工具和三类绑定。
- 陈棋隆：负责 View 层，主要内容包括窗口、画布、输入、HUD 和覆盖层。
- 邱安琪：负责 ViewModel 层，主要内容包括命令回调、状态推进、AI、战斗和胜负判定。

该项目严格完整地按照MVVM架构开发（由于游戏场景中的核心数据（如实体状态、位置等）直接对应视图层的可绘制对象，无需进行二次业务模型转换，因此本架构合并了传统意义上的 Model 层，将其统一内聚于 ViewModel 层中。），在中期报告里有详细说明。
\subsection{游戏运行效果}
相较于中期的第一次迭代版本，我们的最终版本做了界面美化，玩法升级，音效设计。下面分这几个部分展示游戏效果：
\subsubsection{界面显示}
\begin{figure}[H]
% \subfigure 插入多图 计算好间距
\centering    
\subfigure[开始界面] {    
\includegraphics[width=0.45\textwidth]{开始界面.png}  
}     
\subfigure[游戏地图] {     
\includegraphics[width=0.45\textwidth]{游戏地图.png}     
}    
\subfigure[游戏暂停界面] {    
\includegraphics[width=0.45\textwidth]{游戏暂停界面.png}  
}   
\subfigure[boss出现警告] {    
\includegraphics[width=0.45\textwidth]{boss出现警告1.png}  
}
\subfigure[boss出现警告] {    
\includegraphics[width=0.45\textwidth]{boss出现警告2.png}  
}
\subfigure[boss出现] {    
\includegraphics[width=0.45\textwidth]{boss出现.png}  
}
\subfigure[游戏胜利] {    
\includegraphics[width=0.45\textwidth]{游戏胜利.png}  
}

\subsubsection{玩法}
\subsubsubsection{敌人种类}
在游戏中，敌人主要分为冲撞型敌人、伏击类飞行敌人、远程型敌人、远程型敌人2（机器人）和Boss五种类型。每种类型的敌人都有不同的攻击方式和行为模式，增加了游戏的挑战性和趣味性。
\subsection{表格示例}
\begin{table}[H]
  \centering
  \caption{敌人类型}
    \begin{tabular}{|c|c|}
    \hline
    敌人类型 & 特点\\
    \hline
    冲撞型敌人  &   \\
    \hline
    伏击类飞行敌人 &  \\
    \hline
    远程型敌人 &  \\
    \hline
    远程型敌人2（机器人） &  \\
    \hline
    Boss &  \\
    \hline
    \end{tabular}
\end{table}

\begin{figure}[H]
% \subfigure 插入多图 计算好间距
\centering    
\subfigure[冲撞型敌人] {    

\includegraphics[width=0.45\textwidth]{../../code/view/Assets/冲撞型敌人/Attack.png}  
}     
\subfigure[伏击类飞行敌人] {     

\includegraphics[width=0.45\textwidth]{../../code/view/Assets/伏击类飞行敌人/Attack.png}     
}    
\subfigure[远程型敌人] {    

\includegraphics[width=0.45\textwidth]{../../code/view/Assets/远程型敌人/Attack.png}  
}   
\subfigure[远程型敌人2（机器人）] {    

\includegraphics[width=0.45\textwidth]{../../code/view/Assets/远程型敌人2（机器人）/Attack.png}  
}  
\subfigure[Boss] {    

\includegraphics[width=0.45\textwidth]{../../code/view/Assets/Boss/Attack1.png}  
} 
\subfigure[Boss] {    

\includegraphics[width=0.45\textwidth]{../../code/view/Assets/Boss/Attack4.png}  
} 
\caption{敌人类型}  
\end{figure}
在游戏界面上的体现是这样的：
\begin{figure}[H]
% \subfigure 插入多图 计算好间距
\centering    
\subfigure[len-1] {    

\includegraphics[width=0.45\textwidth]{1.png}  
}     
\subfigure[len-2] {     

\includegraphics[width=0.45\textwidth]{2.png}     
}    
\caption{游戏界面体现}  
      
\end{figure}

\subsubsubsection{玩家状态}
玩家可以进行攻击，攻击分为
\begin{figure}[H]
% \subfigure 插入多图 计算好间距
\centering    
\subfigure[玩家被打] {    
\includegraphics[width=0.45\textwidth]{玩家被打.png}  
}     
\subfigure[玩家攻击（脚踢）] {     
\includegraphics[width=0.45\textwidth]{玩家攻击（脚踢）.png}     
}    
\subfigure[玩家攻击（手击）] {    
\includegraphics[width=0.45\textwidth]{玩家攻击（手击）.png}  
}   
\subfigure[玩家空中脚踢] {    
\includegraphics[width=0.45\textwidth]{玩家空中脚踢.png}  
}   

\caption{插入多图}  
      
\end{figure}

\subsubsubsection{补给道具}
在敌人死后会随机掉落血包，拾取后可用于恢复血量/恢复动作耗能
\begin{figure}[H]
% \subfigure 插入多图 计算好间距
\centering    
\subfigure[血包] {    
\includegraphics[width=0.45\textwidth]{敌人死后掉血包.png}  
}     
\subfigure[能量包] {     
\includegraphics[width=0.45\textwidth]{敌人死后掉能量包.png}     
} 

\caption{插入多图}  
      
\end{figure}


\subsubsection{音效}
> 负责的同学补充一下

\subsection{开发工具尝试}
\subsubsection{团队协作与版本管理工具}
为了体现组内协作过程，本项目在开发中持续使用 Git 和 GitHub 记录提交历史。三位成员围绕不同层次并行推进：Common/App 层先确定公共状态和绑定方式，View 层根据公共状态完成输入和绘制，ViewModel 层独立实现规则模拟和状态同步，最后通过主分支合并和冲突处理完成整合。

> 最后加图片

图中的 GitHub 账号与组员对应关系如下：`wys917` 对应苏易文，`qiuanqi70` 对应邱安琪，`odetta1` 对应陈棋隆。从 GitHub Contributors 页面可以看到三位成员均有提交记录，说明项目并不是由单人集中完成，而是按照分层职责持续推进。

需要补充说明的是，上图中统计的提交数和增删行数不能等同于实际工作量，本次任务中我们小组全程保持顺畅的沟通，每个人都很负责的完成了自己的部分，并在有余力时对队友的工作提出建议或帮助，三人的工作量是一致的。

> 最后加图片

上图合并展示了项目的提交历史。可以看到仓库中存在多次围绕 `viewmodel`、`view`、`common` 和 `app` 的独立提交与分支合并，并不是在同一个大文件中串行修改，而是按照 MVVM 层次拆分任务，并持续进行解耦：把公共契约从具体规则中分离出来，让 View 和 ViewModel 可以围绕稳定的 `GameState`、命令回调和通知机制并行开发，最后再通过 merge 将各层成果整合到主分支。

\subsubsection{单元测试}
> 苏易文负责

\subsubsection{持续集成}
> 陈棋隆负责

\section{学习回顾}
> 要提炼核心收获与感悟，邱安琪负责

> 个人总结包括：自身在项目中承担的角色，遇到的主要挑战及应对策略，最终获得的成长与深度反思
\subsection{个人总结：苏易文}

\subsection{个人总结：邱安琪}

本项目中我主要负责 ViewModel 层的设计与实现。在MVVM架构中ViewModel 层位于界面输入和游戏规则之间：一方面接收 View 传来的移动、跳跃、攻击、暂停、确认和重置等语义命令；另一方面驱动游戏模拟，并将规则层的内部状态转换为 View 可以安全读取的公共快照。由于游戏场景中的核心数据（如实体状态、位置等）直接对应视图层的可绘制对象，无需进行二次业务模型转换，因此本架构合并了传统意义上的 Model 层，将其统一内聚于 ViewModel 层中。所以我也负责了游戏模拟的板块。
通过这部分工作，我对如何在 C++ 项目中划分游戏规则、状态同步有了更具体的认识。

\subsubsection{遇到的挑战及应对策略}

\begin{enumerate}
  \item \textbf{遇到的挑战}

    第一项挑战是确定 MVVM 中各层的边界。游戏实体既需要保存规则字段，又需要被绘制成角色、血条和动画，如果直接把内部实体暴露给 View，界面代码就会逐渐依赖 AI、冷却和碰撞细节；如果在 ViewModel 中维护一份完全独立的显示副本，又容易出现规则状态和显示状态不一致。

    第二项挑战是处理连续时间中的离散事件。移动、能量恢复、跳跃高度、攻击命中、敌人受击、投射物和掉落物都依赖时间步，但命中帧、死亡展示结束和 Boss 入场又是离散事件。若没有固定且明确的更新顺序，就可能出现攻击重复命中、掉落物晚一帧生成、Boss 在入场动画完成前被攻击或胜利页面过早出现等问题。

    第三项挑战是让不同敌人的行为既有差异又保持可维护性。巡逻、伏击、冲撞、远程攻击和 Boss 追踪都需要自己的计时器和触发条件，同时又共享受击、死亡、朝向和掉落逻辑。如果把所有逻辑堆在一个函数中，新增敌人或修改某一行为时很容易影响其他类型。

    第四项挑战是处理跨层状态同步。公共 \texttt{GameState} 需要提供 View 绘制所需的尺寸、可见性、HUD、相机和动画状态，但这些信息不应反过来成为规则状态。尤其是死亡实体需要保留一段时间供 View 播放死亡表现，而 Boss 血条只应在 Boss 存活时显示，二者都不能简单地用“是否存在于实体数组”判断。

  \item \textbf{应对策略}

    针对第一项挑战，我把规则实体定义在 \texttt{SimulationTypes.h} 中，把绘制使用的 DTO 放在 Common 层，并通过 \texttt{sync\_state\_from\_simulation()} 集中完成一次映射。View 只能取得 \texttt{const GameState*}，不能修改模拟器内部容器；这样规则层和显示层之间只有明确的单向数据流。

    针对第二项挑战，我把所有逐帧变化集中到 \texttt{GameSimulation::step(float dt)}，并固定更新顺序。攻击使用动作总时长、命中时间和 \texttt{attackHitApplied} 组合实现一次性命中；敌人攻击使用各自的攻击计时器和命中标志；死亡和 Boss 胜利使用独立的展示计时器。对于非法时间步，先用有限性和正数检查拒绝处理，保证模拟状态不会被异常输入破坏。

    针对第三项挑战，我采用“公共前置处理加类型分发”的方式。\texttt{simulate\_ai()} 统一处理冷却、受击恢复和朝向，然后将具体行为交给 \texttt{update\_patroller()}、\texttt{update\_ambusher()}、\texttt{update\_charger()}、\texttt{update\_ranged\_enemy()} 和 \texttt{update\_boss()}。死亡处理和掉落生成集中在 \texttt{process\_enemy\_deaths()}，避免各个 AI 函数重复实现相同的收尾逻辑。

    针对第四项挑战，我区分了“存活”“死亡展示中”和“死亡展示结束”三个概念。死亡实体的 \texttt{alive} 设为 false，但只要 \texttt{deathTimer} 大于零，ViewModel 仍把它作为可见角色放入快照；Boss 则通过 \texttt{boss\_victory\_ready()} 控制最终胜利切换。对于 HUD 和相机等派生数据，ViewModel 在同步阶段统一计算，View 不再自行推导规则结果。
\end{enumerate}

\subsubsection{反思与成长}
这部分工作让我认识到，ViewModel 的价值重要的是建立稳定的输入、状态和通知边界。一个好的 ViewModel 应该让 View 只关心当前应该显示什么、用户发出了什么意图，而不需要知道状态为什么发生变化；同时也应该让规则模拟器可以脱离图形界面独立验证。

我还认识到，游戏规则中的时间顺序本身就是领域知识。攻击命中帧、受击无敌时间、波次过渡、投射物飞行和死亡展示都不是普通的属性赋值，而是状态随时间变化的过程。将这些过程集中到可预测的模拟步中，比在多个回调中分别修改状态更容易推理、调试和测试。

在代码实现上，我从最初关注单个功能是否能够运行，逐渐转向关注不变量是否始终成立。例如，实体列表的第一个元素始终是玩家；已经死亡的敌人不能继续移动或攻击；一次攻击最多结算一次命中；普通波次未完成时不能生成 Boss；无效 \texttt{dt} 不改变模拟状态。这些不变量为后续功能扩展和测试提供了判断标准。

最后，ViewModel 层的实现也让我体会到，模块解耦并不意味着模块之间没有联系，而是要通过清晰的契约进行联系。\texttt{GameSimulation} 负责规则的真实性，\texttt{GameViewModel} 负责规则到界面的翻译，Common 层负责跨层共享的数据契约。

\subsection{个人总结：陈棋隆}

\section{课程改进建议}
> 每个人写一点，我最后再总结一下
> 结合实际体验，提出对课程内容、组织形式或教学方法的具体优化建议，以助力未来学习体验的持续提升

\end{document}
